\section{Notas de aulas}

As aulas foram acompanhadas por definições, observações e problemas destinados
à introdução e à revisão dos conceitos fundamentais de Equações Diferenciais
Ordinárias (EDOs). Nesta seção, são reunidos os principais conteúdos
desenvolvidos ao longo das aulas, incluindo a classificação das EDOs quanto à
ordem, ao grau, à linearidade e às variáveis dependente e independente, bem
como exemplos e problemas utilizados para ilustrar esses conceitos. Também são
apresentados procedimentos de resolução e aplicações das definições estudadas,
servindo como referência para o desenvolvimento dos exercícios e das listas da
disciplina.

\question{
  Considere a equação diferencial:
  \[
    f'(x) + \sin{x} = 0.
  \]
  Determine \(y = f(x)\), sabendo que \(f(0) = 2\).
}

\answer{
  Partindo da equação diferencial,
  \begin{align*}
    f'(x) + \sin{x} &= 0 \\
    f'(x) &= -\sin{x}.
  \end{align*}

  Como \(f'(x) = \dfrac{dy}{dx}\), integra-se ambos os lados em relação a
  \(x\):
  \begin{align*}
    \frac{dy}{dx} &= -\sin{x} \\
    \int dy &= \int -\sin{x}\,dx \\
  \end{align*}

  Dessa forma conseguimos a solução geral:
  \begin{align*}
    y &= \cos{x} + C
  \end{align*}

  Aplicando a condição inicial \(f(0) = 2\),
  \begin{align*}
    f(0) &= \cos{0} + C \\
    2 &= 1 + C \\
    C &= 1
  \end{align*}

  Portanto, a solução particular é:
  \[
    \boxed{y = \cos{x} + 1}.
  \]
}

\newpage

\question{
  Encontre o valor de \(A\) para que
  \[
    y = A\sin{3t}
  \]
  seja solução da equação diferencial
  \[
    \frac{d^{2}y}{dt^{2}} + 2y = 4\sin{3t}
  \]
}

\answer{
  Calculando as derivadas de \(y = A\sin{3t}\), obtém-se
  \begin{align*}
    \frac{dy}{dt} &= 3A\cos{3t} \\
    \frac{d^{2}y}{dt^{2}} &= -9A\sin{3t}
  \end{align*}

  Substituindo \(y\) e sua segunda derivada na equação diferencial,
  \begin{align*}
    -9A\sin{3t} + 2A\sin{3t} &= 4\sin{3t} \\
    -7A\sin{3t} &= 4\sin{3t} \\
    -7A &= 4,
  \end{align*}

  portanto,
  \[
    \boxed{A = -\frac{4}{7}}.
  \]
}

\section*{Ordem, Grau, Linearidade e Variáveis}

\definition{Variável Independente}{
	A \textbf{variável independente} é aquela em função da qual a variável
	dependente é determinada. Em geral, é representada por $x$ ou $t$.

	Por exemplo, na função
	\[
		y=y(x),
	\]
	a variável $x$ é a \textbf{variável independente}.
}

\definition{Variável Dependente}{
	A \textbf{variável dependente} é aquela cujo valor depende da variável
	independente. Em geral, é representada por $y$.

	Por exemplo, na função
	\[
		y=y(x),
	\]
	a variável $y$ é a \textbf{variável dependente}, pois seu valor depende de $x$.
}

\definition{Ordem}{
	A \textbf{ordem} de uma EDO é determinada pela maior derivada da variável
	dependente que aparece na equação.

	Por exemplo,
	\[
		y'''-5xy'+y=e^x
	\]
	é uma EDO de \textbf{terceira ordem}, pois a maior derivada presente é $y'''$.
}

\definition{Grau}{
	O \textbf{grau} de uma EDO é o expoente da maior derivada, desde que a equação
	esteja escrita como um polinômio em relação às derivadas da variável
	dependente.

	Por exemplo,
	\[
		(y'')^3+2y'-y=x
	\]
	é uma EDO de \textbf{grau 3}, pois a maior derivada, $y''$, está elevada à
	terceira potência.
}

\definition{Linearidade}{
	Uma EDO de ordem $n$ é dita \textbf{linear} quando pode ser escrita na forma

	\[
		a_n(x)\frac{d^n y}{dx^n}
		+a_{n-1}(x)\frac{d^{n-1}y}{dx^{n-1}}
		+\cdots
		+a_1(x)\frac{dy}{dx}
		+a_0(x)y
		=g(x),
	\]

	onde $a_0(x),a_1(x),\ldots,a_n(x)$ e $g(x)$ são funções conhecidas da variável
	independente $x$, com $a_n(x)\neq 0$.

	Uma EDO é \textbf{linear} se:

	\begin{enumerate}
		\item Os coeficientes de $y$ e de suas derivadas dependem somente da
		variável independente $x$ ou são constantes;

		\item A variável dependente $y$ e todas as suas derivadas aparecem apenas
		na primeira potência, isto é, com expoente igual a $1$;

		\item Não existem produtos entre $y$ e suas derivadas, nem entre suas
		próprias derivadas.
	\end{enumerate}

	Por exemplo,
	\[
		x^2y''+3xy'-5y=e^x
	\]
	é uma EDO linear, enquanto
	\[
		y''+y^2=0
	\]
	não é linear, pois $y$ aparece elevado ao quadrado.
}

\question{
  Qual função satisfaz a EDO seguinte:
  \[
    y'' + y = 0
  \]
}

\answer{
  A solução geral da EDO
  \[
    y''+y=0
  \]
  é dada por uma combinação linear de $\sin x$ e $\cos x$:

  \[
    \boxed{y(x)=C_1\sin x+C_2\cos x},
  \]

  onde $C_1,C_2\in\mathbb{R}$ são constantes arbitrárias.
}

\section*{Algumas Aplicações Envolvendo EDOs}

\subsection*{Crescimento Populacional: Modelo de Malthus}

\question{
  Quais são os modelos matemáticos que descrevem o crescimento populacional?
}

\subquestion{
  Como podemos considerar os recursos do local para que não ocorra extinção de
  uma espécie?
}

\answer{
  Inicialmente, vamos entender que o problema pode ser modelado por meio de uma
  Equação Diferencial Ordinária (EDO).

	\begin{itemize}
    \item $t$: tempo (variável independente);
    \item $P(t)$: tamanho da população no tempo $t$ (variável dependente);
    \item $\dfrac{dP}{dt}$: variação da população em relação ao tempo.
	\end{itemize}

  Seguimos, agora, para a modelagem da EDO que governa o problema do
  crescimento populacional descrito por Malthus. Perceba que a aplicação diz
  que ``a taxa de variação da população em relação ao tempo é proporcional ao
  tamanho da população presente no tempo $t$''. Portanto,

	\[
		\frac{dP}{dt} = kP(t),
	\]

  onde $k$ é uma constante de proporcionalidade que depende das características
  da população em estudo.

  Neste ponto, já temos a capacidade de reconhecer algumas características da
  EDO desta aplicação:

	\begin{itemize}
		\item É uma EDO de primeira ordem;
		\item O grau é $1$;
		\item É linear.
	\end{itemize}

  De modo muito parecido com o exemplo 1 deste material, vamos aplicar
  separação de variáveis para resolver esta EDO.

	\[
		\frac{dP}{dt}=kP(t)
		\quad\Longrightarrow\quad
		\frac{dP}{P}=k\,dt.
	\]

	Aplicando a integral em ambos os lados, temos:

	\[
		\int\frac{dP}{P}
		=
		\int k\,dt,
	\]

	o que resulta em

	\[
		\ln|P|=kt+c_1.
	\]

  Como queremos obter a função para $P$, precisamos eliminar a função logaritmo
  natural. Isto é feito por meio da função inversa, a exponencial de base $e$.

	\[
		P(t)=e^{kt+c_1}.
	\]

	Utilizando a propriedade da exponencial, obtemos:

	\[
		P(t)=e^{kt}e^{c_1}.
	\]

	Como $e^{c_1}$ é uma constante, podemos definir

	\[
		c=e^{c_1}.
	\]

	Assim, a solução geral da EDO é

	\[
		\boxed{P(t)=ce^{kt}}.
	\]
}

\subsection*{Crescimento Populacional: Modelo Logístico}

\subquestion{
Como podemos considerar os recursos disponíveis no local para que o crescimento
de uma população não seja ilimitado?
}

\answer{
Inicialmente, vamos entender que o problema pode ser modelado por meio de uma
Equação Diferencial Ordinária (EDO).

O modelo logístico considera que o crescimento de uma população não ocorre
indefinidamente. À medida que a população aumenta, os recursos disponíveis
tornam-se mais escassos, fazendo com que a taxa de crescimento diminua.

O modelo é descrito pela EDO

$$
  \frac{dP}{dt}=P(a-bP),
$$

onde:

\begin{itemize}
  \item $t$: tempo (variável independente);
  \item $P(t)$: tamanho da população no tempo $t$ (variável dependente);
  \item $\dfrac{dP}{dt}$: variação da população em relação ao tempo;
  \item $a$: taxa média de crescimento da população, uma constante positiva;
  \item $b$: taxa média de óbitos, uma constante positiva;
  \item $(a-bP)$: termo responsável por representar a redução da taxa de
    crescimento à medida que a população aumenta.
\end{itemize}

Perceba que, diferentemente do modelo de Malthus, a taxa de crescimento agora
depende do tamanho da população. Quando a população é pequena, o termo $(a-bP)$
é positivo e a população tende a crescer. Entretanto, à medida que $P$ aumenta,
esse termo diminui, reduzindo a velocidade de crescimento da população.

Portanto, a EDO

$$
  \frac{dP}{dt}=Pa-bP^2
$$

representa um modelo de crescimento populacional no qual os recursos
disponíveis impõem um limite ao crescimento da população.

\begin{align*}
  \frac{dP}{dt}&=P(a-bP) \\
  \int{\frac{dP}{P(a-bP)}}&=\int{{dt}} \\
  \int{\frac{1}{P(a-bP)}}&=t \\
  \int{\frac{A}{P} + \frac{B}{a-bP}}&=t \\
  \int{\frac{A}{P} + \frac{B}{a-bP}}&=t \\
  aA - bAP + BP&=1 \\
  1&=aA \therefore A=\frac{1}{a} \\
  -bA + B &= 0 \therefore B=\frac{b}{a} \\
  \int{\frac{dP}{P(a-bP)}}&=\int{\frac{\frac{1}{a}}{P}} dP + \frac{\frac{b}{a}}{{a - bP}} = \int{dt} \\
  \frac{1}{a}\int{\frac{1}{P}} + \frac{b}{a} \int{\frac{dP}{a-bP}} &= \int{dt} \\
  \frac{1}{a}ln|{P}| - \frac{1}{a} ln|{a-bP}| &= t + C_1 \\
  \frac{1}{a}[ln|{P}| - ln|{a-bP}|] &= t + C_1 \\
  [ln|{P}| - ln|{a-bP}|] &= a(t + C_1) \\
  ln|\frac{P}{a-bP}| &= a(t + C_1) \\
  \frac{P}{a-bP} &= Ce^{at} \\
  P(1 + bCe^{at}) &= aCe^{at} \\
\end{align*}

Portanto,

\[
  \boxed{P(t) = \frac{aCe^{at}}{1 + bCe^{at}}} \\
\]
}

\section*{EDOs Separáveis}

\definition{EDO Separável}{
Uma Equação Diferencial Ordinária (EDO) é dita \textbf{separável} quando pode
ser escrita de modo que os termos que dependem da variável independente e da
variável dependente possam ser separados.

Em particular, considere uma EDO na forma

\[
	M(t)+N(y)\frac{dy}{dt}=0,
\]

onde $M$ é uma função que depende apenas de $t$ e $N$ é uma função que depende
apenas de $y$.

Multiplicando a equação por $dt$, obtemos

\[
	M(t)\,dt+N(y)\,dy=0.
\]

Podemos, então, separar as variáveis:

\[
	M(t)\,dt=-N(y)\,dy.
\]

Essa forma permite integrar separadamente os termos que dependem de $t$ e de
$y$:

\[
	\int M(t)\,dt+\int N(y)\,dy=0.
\]

Portanto, uma EDO é separável quando pode ser manipulada de modo a obter uma
expressão na qual as variáveis estejam separadas, permitindo a resolução por
meio de integrações.
}

\question{
Determine a solução da EDO

\[
	y'=\frac{1}{y}.
\]
}

\answer{
Inicialmente, fazemos a separação das variáveis. Escrevendo a derivada
como uma razão de diferenciais, temos

\[
	\frac{dy}{dx}=\frac{1}{y}.
\]

Multiplicando ambos os lados por $y\,dx$, obtemos

\[
	y\,dy=x\,dx.
\]

Agora podemos integrar ambos os lados:

\[
	\int y\,dy=\int x\,dx.
\]

Calculando as integrais,

\[
	\boxed{\frac{y^2}{2}
	=
	x+c_1}.
\]
}

\question{
Determine a solução da EDO

\[
  e^yy' - t - tln|t| = e^t\ \ \ \text{Para}\,t>0
\]
}

\answer{
\begin{align*}
  e^y\frac{dy}{dx} - t - tln|t| &= e^t \\
  e^y\frac{dy}{dx} &= t + tln|t| + e^t \\
  \int{e^ydy} &= \int{t dt} + \int{tln|t| dt} + \int{e^t dt} \\
\end{align*}

Portanto,

\[
	\boxed{
    e^y &= \frac{t^2}{2} + \frac{t^2}{2} ln|t| - \frac{t^2}{4} + e^t + C
  }.
\]
}

\question{
Construa a solução geral e determine o domínio do PVI.

\[
  \begin{cases}
    y' = (1 - 2x)y^2 \case \\
    y(0) = \frac{-1}{6} \case
  \end{cases}
\]
}

\answer{
\begin{align*}
  \frac{dy}{dx} &= (1 - 2x)y^2 \\
  \frac{dy}{y^2} &= (1 - 2x)dx \\
  \int{y^{-2}dy} &= \int{(1 - 2x)dx} \\
  \frac{y^{-1}}{-1} &= x - x^2 + C \\
  \frac{1}{y} &= -x + x^2 - C \\
\end{align*}

Portanto,

\[
	\boxed{
  y &= \frac{1}{x^2 - x - C}
  }.
\]

Aplicando $y(0)$,

\begin{align*}
  \frac{-1}{6} &= \frac{1}{- C} \therefore C=6
\end{align*}

Sendo assim, o domínio de $Y_1$ é dado por

\begin{align*}
  x^2-x-6 &\neq 0 \\
  (x-3)(x+2) &\neq 0 \\
  x &\neq -2,\,3.
\end{align*}

Portanto,

\[
  D_y = \mathbb{R} \setminus \{-2,3\}.
\]
}
